% NTCIR-19 FEHU Task 1 working preview.
% Prose source: paper_draft.md
% This file follows the official NTCIR-19 acmart setup and adds the a4paper
% option at the user's request. Placeholder space is deliberate and forms part
% of the eight-page budget check.
\documentclass[sigconf,article,a4paper]{acmart}

\settopmatter{printacmref=false}
\renewcommand\footnotetextcopyrightpermission[1]{}
\pagenumbering{gobble}

\usepackage{amsmath}
\usepackage{array}
\usepackage{booktabs}
\usepackage{float}
\usepackage{tabularx}
\usepackage{xcolor}
\geometry{a4paper}
\setlength{\emergencystretch}{1em}

\definecolor{placeholderbg}{RGB}{247,247,247}
\definecolor{placeholderline}{RGB}{185,185,185}
\newcommand{\placeholdernote}[1]{%
  \par\smallskip
  \noindent\fcolorbox{placeholderline}{placeholderbg}{%
    \parbox{0.94\columnwidth}{\centering\small\color{gray}#1}}
  \par\smallskip
}
\newcommand{\placeholderblock}[2]{%
  \par\smallskip
  \noindent\fcolorbox{placeholderline}{placeholderbg}{%
    \parbox[c][#2][c]{0.94\columnwidth}{%
      \centering\small\color{gray}#1}}
  \par\smallskip
}

\title{WAIPS at NTCIR-19 FEHU Task 1: Label-Aware Prompt Ensembling and Gated Fusion for Human-Value Recognition}

% TODO(camera-ready): confirm author order, affiliations, countries, and e-mail.
\author{Peiyang Huang}
\affiliation{%
  \institution{Affiliation to be confirmed}
  \country{China}}
\author{Dongyang Liu}
\affiliation{%
  \institution{Affiliation to be confirmed}
  \country{China}}
\author{WAIPS Team}
\affiliation{%
  \institution{Additional authors to be confirmed}
  \country{China}}

\begin{abstract}
We present the three WAIPS runs submitted to NTCIR-19 FEHU Task 1, which asks
systems to recognize actor-conditioned human values in news and, for
fine-grained values, predict direction. Our development experiments exposed a
recurring trade-off. Conservative prompt ensembles preserve direction but miss
infrequent values; recall-oriented methods recover more labels but add false
positives. WAIPS-01 selectively
adds labels from a targeted prompt to a conservative majority according to
label-wise analysis on development data. WAIPS-03 augments difficult directed
labels and tunes DeBERTa-v3-base with LoRA and masked partial-label supervision. The encoder
probabilities join heterogeneous prompt jurors in calibrated, frequency-gated
veto and addition rules, followed by focused direction verification. WAIPS-03
scores 0.5808 Micro-F1 and 0.3952 Macro-F1 on official Task 1b, the highest
Macro-F1 among participant runs. WAIPS-01 has the lowest participant
Direction Reverse Rate (DRR), 0.0345. Targeted augmentation
improves score separation for 18 low-frequency labels, and E6b supplies direct
tail recall. Their combination under conservative gates supports macro
performance, although the missing final no-DeBERTa ablation prevents assigning
the official gain to one component.
\end{abstract}

\keywords{human values, multi-label classification, long-tail learning, data
augmentation, prompt ensemble, calibrated fusion}

\begin{document}
\maketitle
\pagestyle{plain}

\section*{Team Name}
WAIPS

\section*{Subtasks}
Task 1a and Task 1b

\section{Introduction}

Human values are often expressed indirectly in news. A report may describe the
same decision as protecting public safety for one actor and restricting freedom
for another. Recognizing a value therefore requires more than locating a
topical phrase: the system must identify whose action is being judged, which
value it bears on, and whether the action aligns with or contradicts that
value. The FEHU Task 1 makes these choices explicit. Task 1a asks for
directionless Level-2 values for each article--actor pair, while Task 1b
evaluates directed Level-1 values. We refer readers to the Task Overview by
Wang et al. for the collection construction, taxonomy, and official evaluation
protocol.

Three properties of the task shaped our work. First, the output space is
strongly imbalanced. A small group of values appears throughout the training
set, while many directed Level-1 classes have only a few examples. A system can
therefore obtain reasonable Micro-F1 while never predicting several valid
labels. Second, increasing recall is unusually expensive: every extra value
creates another opportunity for both a false positive and a direction
reversal. Third, the available models make different errors. Prompt variants
disagree in useful ways, and the encoder can expose weak tail-label evidence
that is not reliable enough to use as a final decision by itself.

Rather than forcing these components into a single voting rule, we assigned
them different jobs. WAIPS-01, developed by Liu, starts with a conservative
prompt majority and selectively restores predictions for labels on which a
targeted prompt was demonstrably better. A separate augmentation track
diagnoses difficult directed labels, generates focused training cases, and
trains a parameter-efficient DeBERTa classifier to supply continuous evidence.
WAIPS-03 (E7-G4) uses E6b, a recall-oriented prompt system with
self-segmentation, as its primary generator; other prompt runs and the
augmented encoder act as jurors.
A calibrated gate removes weak frequent-class predictions and admits a small
number of supported additions. Direction is then checked only when jurors
explicitly disagree.

This division of labour is also how we interpret the results. E6b supplies much
of the direct tail recall, the DeBERTa track improves representation of
selected low-frequency labels and contributes an auxiliary probability signal,
and the fusion layer controls the precision cost of combining them. We report
both F1 and DRR throughout because the official runs occupy visibly different
points on that trade-off.

\section{Task Setting and Submitted Runs}

For each actor listed in an article, Task 1b requires a set of (Level-1 value,
direction) pairs. We treat the two directions of each of 54 Level-1 values as
108 atomic classes. Task 1a is obtained by mapping the predicted Level-1 values
to their Level-2 parents and discarding direction. Development experiments use
the official 177-article development set. The frozen inference pipelines were
then applied to 175 test articles containing 479 released article--actor pairs;
no test labels or Task-2 annotations were used.

We follow the official evaluator and the Task Overview in treating Micro-F1 as
the primary ranking measure. Macro-F1 is averaged over labels that occur at
least once in the test gold data, making it a useful but support-dependent view
of tail behavior. DRR counts opposite-direction predictions for gold
actor--value pairs with a unique direction. A missed value does not count as a
reversal, so a low DRR must be read together with F1 rather than as a standalone
accuracy measure.

The three submissions were deliberately different. WAIPS-01 is Liu's
label-aware prompt ensemble, described in Section~\ref{sec:liu}. WAIPS-02 uses
a three-way majority vote over DeepSeek Pro prompt variants and uses a Flash
prediction, when available, to arbitrate direction. WAIPS-03 is the E7-G4
heterogeneous fusion described in Sections~\ref{sec:e7} and~\ref{sec:results}.
This paper focuses on WAIPS-01 and WAIPS-03 because they embody the two main
methodological lines; WAIPS-02 is retained as a useful prompt-only reference
point.

\begin{table}[H]
  \caption{The three submitted WAIPS runs.}
  \label{tab:runs}
  \small
  \begin{tabularx}{\columnwidth}{@{}lX@{}}
    \toprule
    Run & Short description \\
    \midrule
    WAIPS-01 & Label-aware prompt ensemble; direction-stable selective recovery. \\
    WAIPS-02 & Three-way prompt majority with direction arbitration. \\
    WAIPS-03 & E7-G4 gated heterogeneous fusion. \\
    \bottomrule
  \end{tabularx}
\end{table}

WAIPS-02 provides a useful middle point between the other two submissions. One
DeepSeek Pro prompt is the baseline, one is compact and strict, and one uses
targeted balanced few-shot examples; all run at temperature zero. A label is retained
when at least two runs emit the same actor, value, and direction.
The Flash run is then consulted only for direction: when it predicts the same
actor and value with the opposite direction, the majority direction is
replaced. Task 1a is deterministically projected from the resulting Task-1b
set. Unlike E7-G4, this run has no learned calibrator, encoder features,
frequency gate, or per-cell QA call.

\section{Label-Aware Prompt Ensembling (WAIPS-01)}
\label{sec:liu}

% The following blank area reserves approximately 1.7 A4 pages in two-column
% layout. Each marker is intentionally brief; Liu's prose should replace the
% entire reserved area, not be appended after it.
\placeholdernote{Reserved for Dongyang Liu (WAIPS-01), part 1 of 5.}
\vfill\newpage
\placeholdernote{Reserved for Dongyang Liu (WAIPS-01), part 2 of 5.}
\vfill\newpage
\placeholdernote{Reserved for Dongyang Liu (WAIPS-01), part 3 of 5.}
\vfill\newpage
\placeholdernote{Reserved for Dongyang Liu (WAIPS-01), part 4 of 5.}
\vfill\newpage
\placeholderblock{Reserved for Dongyang Liu (WAIPS-01), part 5 of 5.}{0.31\textheight}

\section{Diagnosis-Guided Augmentation for Rare-Label Evidence}
\label{sec:augmentation}

\subsection{From label diagnosis to generation targets}

Our initial DeBERTa classifier used a 108-dimensional sigmoid head and a single
global threshold. Per-label analysis showed that this model had effectively
learned no usable ordering for many low-frequency directed labels. We therefore
selected targets at the directed Level-1 class evaluated by Task 1b, rather
than augmenting undirected values uniformly. A class was eligible when its
official training count was at most 130 and its
development support was at least three. Ranking the eligible classes by
diagnostic evidence produced 18 frozen targets: 14 contradictory and four
aligned classes, spanning 12 Level-2 groups.

The data were constructed in two rounds. Round 1 covered six Level-1 values in
both directions. DeepSeek V4 Pro generated 422 candidates; deterministic checks
and an independent LLM semantic screen retained 340 silver records. These
records were machine-screened, not manually verified. Round 2 concentrated on
the 18 directed targets and allocated exactly 50 records to each target. Its
900 records mix targeted in-context generation, direction contrast cases,
situation transformations, and hard negatives (Table~\ref{tab:generation}).
The hard-negative group includes 36
explicitly known negative labels. All Round-2 records passed schema, taxonomy,
actor mention, direction, and duplicate checks, but they did not receive
semantic LLM screening or human review. We keep this provenance distinction
because it limits what can be claimed about data quality.

\begin{table}[H]
  \caption{Round-2 generation strategies.}
  \label{tab:generation}
  \scriptsize
  \begin{tabularx}{\columnwidth}{@{}XrX@{}}
    \toprule
    Strategy & $n$ & Supervision and purpose \\
    \midrule
    Targeted ICL & 450 & Full multi-label; varied contexts and co-values. \\
    Direction contrast & 180 & Target-only; 90 aligned/contradictory pairs. \\
    Transformation & 180 & Target-only; different actor, action, or setting. \\
    Hard negative & 90 & Target-only; 36 explicit negative labels. \\
    \bottomrule
  \end{tabularx}
\end{table}

Generation was deliberately more constrained than free-form paraphrasing.
Each request fixed the actor, target identity, direction, and a news-style
evidence requirement. Local validators rejected malformed mappings, absent
actors, unsupported evidence sentences, forbidden meta-language, and
duplicates before a record could enter a training input. The generation gates
rejected and regenerated 67 Round-1 and 472 Round-2 candidates; the later
Round-1 semantic screen separately excluded 82 of 422 records. The distinction
between structural gates and semantic review is important: the former
establishes format and internal consistency, not that every synthetic label is
correct.

\subsection{Partial-label training}

The final training set contains 1,228 official articles, 340 Round-1 records,
and 900 Round-2 records: 2,468 articles and 4,575 actor examples in total.
Synthetic news vignettes are often only partially annotated. Treating every
unlisted class as negative would create hundreds of false-negative training
signals, especially for the single-target records. We therefore mark 3,785
examples as fully supervised and 790 as target-only. For a target-only
example, binary cross entropy is evaluated only on its known positive entries
and any explicitly known negatives; the remaining label positions receive a
zero loss mask.

We fine-tune DeBERTa-v3-base with LoRA adapters on the query
and value projections ($r=16$, $\alpha=32$, dropout 0.1) and a 108-way
classification head. Inputs pair the article text with the actor name, use a
maximum length of 512, and are optimized with class-weighted binary cross
entropy. The positive weight for each class is
$\sqrt{n_{\mathrm{negative}}/n_{\mathrm{positive}}}$, clipped to $[1,10]$ and
computed over supervised entries only. Models are selected by development
Task-1b Micro-F1 and decoded at a global threshold of 0.30. In E7-G4 the
encoder never assigns the final direction: we export all 108 sigmoid
probabilities and use them only as continuous jury features.

For example $i$ and directed class $k$, let $m_{ik}$ be one for an observed
label decision and zero for an unknown entry. Training minimizes
\begin{equation}
\mathcal{L}=\frac{\sum_{i,k}m_{ik}w_k\operatorname{BCE}(y_{ik},z_{ik})}
{\sum_{i,k}m_{ik}}.
\end{equation}
All mask entries are one for official and fully labelled records. On
target-only records, only the intended positives and explicitly supplied hard
negatives contribute. This does not rebalance the dataset to 50:50; it prevents
missing synthetic annotations from becoming certain negatives. The distinction
matters in a 108-class problem, where a single partially annotated example
would otherwise create more than one hundred negative loss terms.

\subsection{What the augmentation did---and did not---establish}

The clearest comparison is score separation on the 18 targets.
With official data only, mean target AUC was 0.536 and mean per-target oracle
F1 was 0.053. The augmented v2 checkpoint raised these values to 0.697 and
0.112; AUC improved for 14 of 18 targets and oracle F1 for 16 of 18. Thus the
generated supervision changed the ranking of target positives, rather than
merely shifting an unchanged distribution across the 0.30 threshold.

\begin{table}[t]
  \caption{Score-level tail diagnostics on development data.}
  \label{tab:tail}
  \small
  \begin{tabular}{@{}lrr@{}}
    \toprule
    Diagnostic & Official & Augmented v2 \\
    \midrule
    Mean AUC (18 targets) & .536 & \textbf{.697} \\
    Mean oracle F1 & .053 & \textbf{.112} \\
    AUC $\leq .55$ & 11/18 & \textbf{3/18} \\
    \bottomrule
  \end{tabular}
\end{table}

The three-seed comparison isolates the 900 Round-2 records. It gives a second
view of the same effect: the official-only condition remains close to chance
(mean target AUC 0.513), whereas the 900 records raise AUC to 0.624 with
assume negative supervision and 0.577 with masking. Both improvements have the
same sign for all three matched seeds. With the complete 340+900 data, masked
supervision was retained because it had the better mean full-data result and
protected the sparsely annotated Round-1 records; its mean target AUC was
0.632, but with a standard deviation of 0.105.

\begin{table}[t]
  \caption{Three-seed target analysis (mean $\pm$ SD).}
  \label{tab:seeds}
  \small
  \begin{tabular}{@{}lrr@{}}
    \toprule
    Training condition & AUC & Oracle F1 \\
    \midrule
    Official only & $.513\pm.030$ & $.048\pm.010$ \\
    +900, assume negative & $\mathbf{.624\pm.011}$ & $\mathbf{.095\pm.014}$ \\
    +900, masked & $.577\pm.028$ & $.080\pm.019$ \\
    +340+900, masked & $.632\pm.105$ & $.111\pm.049$ \\
    \bottomrule
  \end{tabular}
\end{table}

The gain is useful but not a claim of a solved tail classifier. Three
full-data seeds showed substantial variance: target true positives were 0, 0,
and 37, and overall Macro-F1 ranged from 0.022 to 0.182. This is why the
encoder is an auxiliary source rather than a submitted standalone model. It is
also why we describe the final macro result as a joint effect of tail-oriented
generation, E6b's direct recall, and gated fusion. A clean final-fusion
ablation without DeBERTa is not available, so the official E7-G4 gain cannot be
assigned causally to augmentation alone.

\section{E7-G4: Frequency-Gated Heterogeneous Fusion}
\label{sec:e7}

\subsection{Members with deliberately different roles}

E7-G4 starts from E6b, a prompt pipeline that first segments an article into
six to nine coherent developments and then predicts values for each actor using
a fixed few-shot scaffold. This granularity constraint is small but
consequential. Its predecessor E6 produced a median of 12 segments and achieved
high recall with many false positives. Asking the segmenter to merge closely
related sentences reduced the median to eight. On development data, E6b
improved over E6 from 0.5412 to 0.5496 Micro-F1 and from 0.4079 to 0.4124
Macro-F1. Relative to the earlier E1 prompt run, its Micro-F1 changed only
slightly (0.5480 to 0.5496), but Macro-F1 rose by 0.0405. We therefore use E6b
as the primary generator: it already carries the broad coverage that a later
fusion stage cannot manufacture.

The jurors serve different roles. E1 uses a balanced article-level prompt. E6
instead uses finer self-segmentation to favor recall. E6c runs DeepSeek Pro
without thinking over the E6b segments. It has the strongest individual
Task-1a development result (0.6774) but weak direction behavior; we therefore
use it as value evidence, not as a direction authority. The local Liu component is one
targeted balanced few-shot prompt run and is distinct from the
complete label-aware WAIPS-01 method in Section~\ref{sec:liu}. Finally, the
augmented DeBERTa model contributes continuous probabilities rather than
discrete predictions.

\begin{table}[t]
  \caption{Development profiles and assigned E7 roles.}
  \label{tab:members}
  \scriptsize
  \begin{tabularx}{\columnwidth}{@{}lrrX@{}}
    \toprule
    Member & Micro/Macro & DRR & E7 role \\
    \midrule
    E1 & .5480/.3719 & .0347 & Balanced juror \\
    E6 & .5412/.4079 & .0504 & Recall juror \\
    \textbf{E6b} & \textbf{.5496/.4124} & .0459 & Primary generator \\
    E6c & .5343/.3725 & .0657 & Value/disagreement juror \\
    Liu local & .5347/.3385 & \textbf{.0230} & Precision juror \\
    DeBERTa & --- & --- & Continuous evidence \\
    \bottomrule
  \end{tabularx}
\end{table}

\begin{figure}[t]
  \centering
  \fcolorbox{placeholderline}{white}{%
    \begin{minipage}{0.91\columnwidth}
      \centering\small
      Article + actor\\[2pt]
      $\downarrow$\\[-1pt]
      E6b primary predictions\\[2pt]
      $\searrow$\hspace{22pt}$\swarrow$\\[-1pt]
      \begin{tabular}{c@{\hspace{10pt}}c}
        E1 / E6 / E6c / Liu & DeBERTa probabilities
      \end{tabular}\\[2pt]
      $\downarrow$\\[-1pt]
      Calibrated veto + capped additions (G3)\\[2pt]
      $\downarrow$\\[-1pt]
      Disagreement-triggered direction QA (G4)\\[2pt]
      $\downarrow$\\[-1pt]
      Task-1b $\rightarrow$ deterministic Task-1a projection
    \end{minipage}}
  \caption{E7-G4 separates value generation, calibrated fusion, and direction repair.}
  \Description{The E6b primary predictions, four prompt jurors, and DeBERTa
  probabilities feed calibrated veto and addition rules, followed by focused
  direction quality assurance and deterministic Task-1a projection.}
  \label{fig:e7}
\end{figure}

\subsection{Calibrated veto and additions}

For an actor and directed value cell $c$, the removal score combines
exact-direction votes, value-only votes, and the encoder's directionless value
score:
\begin{equation}
s_{\mathrm{rm}}(c)=n_{\mathrm{exact}}(c)+0.25n_{\mathrm{value}}(c)
 +p_D(\mathrm{value}(c)).
\end{equation}
For a candidate not emitted by E6b, the addition score is stricter:
\begin{equation}
s_{\mathrm{add}}(c)=n_{\mathrm{exact}}(c)+p_D(c).
\end{equation}
Pool-adjacent-violators calibration maps each monotone score to an empirical
probability of being gold. During development, calibration is two-fold
cross-fitted by article, so an article's own labels do not determine its
mapping. For test inference, the mapping is fitted once on the full development
set and applied unchanged. The decision threshold follows the plug-in
expected-F1 rule, $\tau=F1_{\mathrm{primary}}/2$; the frozen full-development
value is 0.27479.

The rules are intentionally asymmetric. A primary cell is removed only when
its calibrated probability is below $\tau$, it belongs to one of 38 frequent
directed classes (official training count greater than 130), and the actor
retains at least one label. A non-primary cell can be added only when it was
proposed by a juror, its calibrated probability reaches $\tau$, it is in the
same frequent tier, and E6b does not already hold the value in either direction.
Additions are capped at two per actor. The frequency gate matters: removing
weak tail predictions reduced Macro-F1 sharply in the ungated development
control. It also means that the 18 augmentation targets are not directly
switched on or off by the G3 actions; the encoder supplies broader score
evidence, while E6b remains the main direct source of tail labels.

\subsection{Focused direction repair}

The last stage addresses a different error. A fused prediction is sent to
direction QA only if at least one juror predicts the same value in the opposite
direction and the fused set does not already contain both directions. DeepSeek
Flash then receives the article, actor, and one value and returns only
\texttt{aligned} or \texttt{contradictory}. On development data, 176 cells were
flagged and 46 were flipped; seven failed all retries and were retained with
explicit failure records. On test, 207 cells were flagged, 77 were flipped,
and five failed and were retained. Because this stage never introduces a new
value, it can repair direction without reopening the value-selection problem.

\subsection{Frozen development-to-test transfer}

The complete test path was executed as a fit--apply procedure. All prompt
members read only the official Task-1 test article file; E6 and E6b generated
their own segmentations because official test subevents were unavailable. The
removal and addition calibrators were fitted on all development data and loaded
unchanged for test. We also froze the frequency tier, score weights, threshold,
maximum of two additions, and QA trigger. Every empty
API response or parsing failure was retried and then either caused the stage to
fail or, for the predeclared QA behavior, was logged while retaining the
original direction.

Before submission, both files were checked for complete article coverage,
valid actor membership, taxonomy ranges, duplicate tuples, and deterministic
Task-1a projection. The final WAIPS-03 output contained 2,764 directed labels,
covering 93 of the 108 classes and all 54 Level-1 values. These counts cannot
substitute for test labels, but the removal/addition profile was close to
development and gave no operational indication of a broken transfer pipeline.

\section{Experiments and Results}
\label{sec:results}

\subsection{Development results}

Table~\ref{tab:dev} follows the construction of E7-G4. The first comparison is
important for the paper's macro-F1 interpretation: E6b already raises
Macro-F1 by 0.0405 over E1. The calibrated veto (G1) mainly restores precision.
Re-enabling supported additions (G3) provides the largest single fusion
increment, and direction QA (G4) produces a smaller but consistent improvement
in Micro-F1, Macro-F1, and DRR.

\begin{table}[H]
  \caption{Development results. T18 is Micro-F1 over the 18 augmentation targets.}
  \label{tab:dev}
  \centering
  \resizebox{\columnwidth}{!}{%
  \begin{tabular}{@{}lrrrrrr@{}}
    \toprule
    System & 1b Micro & Precision / Recall & 1b Macro & 1a Micro & DRR & T18 \\
    \midrule
    E1 & .5480 & .541 / .556 & .3719 & .6574 & .0347 & .2460 \\
    E6b primary & .5496 & .502 / .608 & .4124 & .6621 & .0459 & .3145 \\
    G1: calibrated veto & .5608 & .559 / .563 & .4146 & .6587 & \textbf{.0374} & .3145 \\
    G3: veto + additions & .5790 & .539 / .626 & .4266 & .6874 & .0450 & .3145 \\
    \textbf{G4: G3 + direction QA} & \textbf{.5817} & \textbf{.542 / .629} & \textbf{.4308} & \textbf{.6874} & .0419 & \textbf{.3264} \\
    \bottomrule
  \end{tabular}}
\end{table}

The G4--E6b Micro-F1 difference is 0.0321. In a paired article bootstrap with
10,000 resamples, the 95\% interval runs from 0.0209 to 0.0433. This interval
measures article-sampling variation; it does not measure nondeterminism across
LLM calls. Operationally, G1 removed 484 predictions on development data and
G3 admitted 364. The test pipeline was frozen before inference and produced a
similar action profile: 529 removals and 365 additions. These counts are useful
sanity checks, not test-set quality estimates.

Two controls clarify why the final rule is gated. Removing Liu from the veto
jury changes Micro-F1 only from 0.5608 to 0.5619, so Liu is not load-bearing
for removals; its precision profile is more useful when additions are
considered. Conversely, allowing vetoes on all classes gives 0.5703 Micro-F1
but only 0.3900 Macro-F1. This is the failure mode the frequency gate was
designed to prevent: a weak auxiliary score can look useful in aggregate by
suppressing difficult labels that macro-F1 is intended to expose.

\subsection{Where the macro gain enters}

The development path makes the macro story more specific than the final run
comparison alone. Moving from E1 to E6b raises Macro-F1 from 0.3719 to 0.4124
(+0.0405) while changing Micro-F1 by only +0.0016. The frequent-class veto
then adds +0.0022 Macro-F1, supported additions add a further +0.0120, and
direction QA adds +0.0042. The 18-target Micro-F1 is 0.3145 for E6b and remains
unchanged through G1 and G3; direction QA raises it to 0.3264.

This decomposition matches the assigned roles. E6b is the component that
directly emits most of the additional tail labels. Augmentation makes the
DeBERTa score space less blind to rare labels and contributes a continuous
confidence signal, but the frequent-class gate prevents G3 from directly
adding the 18 selected targets. Its support for the final macro-oriented system
is therefore indirect and shared with the prompt jury. We regard the official
macro advantage as consistent with the combined design, not as a demonstration
that augmented DeBERTa alone caused it.

\subsection{Official results}

The Task Overview reports four evaluated run packages from two teams: three
from WAIPS and one from IIUoT. The results show different trade-offs rather than
one uniformly dominant run. WAIPS-02 gives our highest Task-1a Micro-F1,
WAIPS-03 gives our strongest Task-1b Micro- and Macro-F1, and WAIPS-01 gives the
lowest DRR among participant runs. Across teams, IIUoT-01 leads the primary
Micro-F1 ranking on Tasks 1a and 1b, while WAIPS-03 has the highest participant
Task-1b Macro-F1. The organizer baseline is an unranked reference and exceeds
every participant run on Task-1b Micro- and Macro-F1. We therefore state the
macro result as a participant comparison, not as overall task dominance.

\begin{table}[H]
  \caption{Official Task-1 results. Each F1 entry is Micro / Macro. The
  organizer baseline is unranked; bold marks the best participant.}
  \label{tab:official}
  \centering
  \scriptsize
  \resizebox{\columnwidth}{!}{%
  \begin{tabular}{@{}lrrr@{}}
    \toprule
    System / run & Task 1a F1 & Task 1b F1 & DRR \\
    \midrule
    Organizer baseline & .6782 / .5405 & .6284 / .4402 & .0386 \\
    IIUoT-01 & \textbf{.7042 / .5808} & \textbf{.5979} / .3930 & .0613 \\
    WAIPS-01 & .6589 / .5366 & .5543 / .3277 & \textbf{.0345} \\
    WAIPS-02 & .6930 / .5646 & .5496 / .3580 & .0395 \\
    \textbf{WAIPS-03} & .6862 / .5705 & .5808 / \textbf{.3952} & .0477 \\
    \bottomrule
  \end{tabular}}
\end{table}

The development and official results tell a consistent, if not perfectly tidy,
story. Broad value recovery helps macro-F1, but it increases exposure to
direction mistakes. WAIPS-01 is the conservative end of this spectrum. E7-G4
accepts a higher DRR in exchange for more labels and higher Task-1b F1. Its
targeted direction QA reduces the development reversal rate after additions,
but does not erase the underlying coverage--direction trade-off on the official
test set.

\subsection{Error patterns and limitations}

Two errors remain common. The first is over-generation of plausible values for
a nearby actor or for the article in general. Self-segmentation makes more local
evidence visible, but a value can still be attached to an actor who is mentioned
rather than responsible. The calibrated veto repairs some of these errors only
in frequent classes. The second is direction ambiguity. Jurors often agree that
a value is relevant while disagreeing about whether the actor supports or
undermines it; focused QA helps, but ambiguous reporting and multiple actions
by one actor remain difficult.

There are also limits to the evidence behind our interpretation. The fusion
weights, frequency gate, and threshold family were designed through
development-era analysis, even though calibration was cross-fitted and the test
pipeline was frozen. The Round-2 synthetic data were structurally checked but
not semantically screened, and masked DeBERTa training has high seed variance.
Most importantly, the final fusion was not rerun in a controlled no-DeBERTa
condition, so the augmentation track's causal share of the official macro
result is unknown. We can safely say that it improves low-frequency score
separation and provides auxiliary evidence; E6b and the LLM jury remain
essential. Finally, results from two participating teams should not be read as
a broad leaderboard conclusion.

\section{Related Work}
\placeholderblock{Reserved for concise related work (prompt ensembles,
long-tail multi-label learning, targeted augmentation, DeBERTa/LoRA,
calibration, and human-value recognition).}{0.22\textheight}

\section{Conclusion}

Our Task-1 systems were built by treating label recovery, precision control,
and direction assignment as separate problems. Liu's WAIPS-01 method uses
label-level diagnosis to recover selected misses without abandoning a
conservative prompt majority. The augmentation track gives a LoRA-tuned
DeBERTa encoder substantially better score-level separation on selected rare
directed labels, while E6b supplies strong direct tail recall through
granularity-controlled self-segmentation. E7-G4 combines these heterogeneous
signals with calibrated, frequency-gated veto and addition rules and repairs
direction only where jurors disagree. The official results reflect the
intended trade-off: WAIPS-03 reaches 0.5808 Micro-F1 and 0.3952 Macro-F1 on
Task 1b, while WAIPS-01 remains more direction-stable. A clean component
ablation and more stable tail training are the clearest next steps.

\section*{References}
\placeholderblock{Reserved for the FEHU Task Overview, taxonomy/resource,
DeBERTa-v3, LoRA, prompt ensemble, long-tail learning, synthetic augmentation,
and calibration references.}{0.28\textheight}

\end{document}
